% =====================================================================
% Cover letter VARIANT A (phenomenon-led) — NMR in Biomedicine
% EIC: Professor John R. Griffiths
% Built with tectonic: `tectonic cover_letter_phenomenon.tex`
% =====================================================================
\documentclass[11pt]{letter}
\usepackage[margin=1in]{geometry}
\usepackage{microtype}
\usepackage[hidelinks]{hyperref}

\signature{Avery Karlin, BA \\ Department of Applied Physics and Applied Mathematics,
Columbia University \\ Department of Computer Science, University of Colorado Boulder \\
\texttt{ak5232@columbia.edu}}

\address{Avery Karlin \\ Columbia University \\ New York, NY, USA}

\begin{document}

\begin{letter}{Professor John R.\ Griffiths \\ Editor-in-Chief, \textit{NMR in Biomedicine}}

\opening{Dear Professor Griffiths,}

I am pleased to submit the manuscript ``Pseudo-diffusion is weakly identifiable in
practice: a reproducible boundary-railing signature of $D^{*}$ in open in-vivo IVIM, and
the shape-aware intervals that resolve it'' for consideration in \textit{NMR in
Biomedicine}.

The work begins from a concrete, real-data observation. Pseudo-diffusion ($D^{*}$) is the
least identifiable parameter of the intravoxel incoherent motion (IVIM) model, a fact
long known in principle from simulation. We show it has a direct, measurable fingerprint
in real acquisitions: on the open OSIPI in-vivo abdomen scan, a conventional
box-constrained least-squares fit drives $D^{*}$ onto a parameter bound --- it
``rails'' --- in 54.2\% (95\% CI 52.0--56.4) of high-SNR voxels. This is not an artefact
of one analysis pipeline: an independent, clean-room reimplementation reproduces a
previously reported 54.7\% without using its code, and the phenomenon generalises across
selection, scanner, and organ --- 47.8\% over the full abdominal ROI and 43.7--73.4\% in
an independent multi-subject liver DWI cohort, the sparser b-value scheme railing more,
exactly as identifiability predicts. Railing is an assumption-free read-out of weak
identifiability: it needs no ground truth and no trust in a noise model, only the
optimiser's own output.

We then trace what this does to the uncertainty estimates a clinician would act on, and
how to repair it. The calibration failure is \emph{conditional} --- symmetric Gaussian
intervals under-cover $D^{*}$ specifically in the high-$D^{*}$ regime (central-95\%
coverage 0.63) while remaining near-nominal elsewhere --- and is resolved by shape-aware
posterior quantile intervals and an amortized neural posterior, up to an irreducible
high-$D^{*}$ identifiability wall. Reporting coverage conditionally, and under an
explicitly stated Cram\'er--Rao convention, is what makes the claim honest and
reproducible rather than a single fragile headline.

The manuscript fits the scope of \textit{NMR in Biomedicine} in both its principal
categories: it is (a) an original research study of an in-vivo IVIM phenomenon on open
biomedical MR data, and (b) a methodological contribution to how quantitative-MR
uncertainty should be characterised and reported. Every reported number traces to a
single seeded reproduction that regenerates with one command, and the primary phenomenon
rests on openly licensed in-vivo data (OSIPI TF2.4, Zenodo 14605039; TCGA-LIHC via The
Cancer Imaging Archive), so a reader can check it independently. To our knowledge this is
the first systematic real-data quantification of $D^{*}$ boundary-railing across cohorts,
paired with an honest conditional-coverage characterisation and its resolution; it is
distinct from, and complementary to, learned-uncertainty frameworks for IVIM such as
Casali et al.\ (2026, this journal), whose object is a particular network's internal
uncertainty budget rather than an estimator-agnostic real-data signature.

The manuscript is original, has not been published previously, and is not under
consideration elsewhere. The author declares no competing interests and takes full
responsibility for the integrity of the work; a declaration of AI-assisted writing is
included. I would be glad to suggest suitable reviewers and to answer any questions.

\closing{Sincerely,}

\end{letter}
\end{document}
